% SPDX-License-Identifier: GPL-3.0-or-later
% Copyright (C) 2026 Alessandro Greco
%% Русское альтернативное описание EPUB.
\altfig{Телескопическое построение цепи Tor из трёх переходов: клиент, R1 (сторожевой), R2 (промежуточный), R3 (выходной) показаны отдельными колонками. Клиент обменивается CREATE2/CREATED2 с R1 и выводит ключи K1. Затем EXTEND2 идёт внутри существующей цепи: R1 обрабатывает запрос и отправляет CREATE2 к R2, а ответы CREATED2/EXTENDED2 возвращают данные для вывода K2. Для третьего перехода EXTEND2 проходит через сторожевой узел к R2; R2 обрабатывает его и посылает CREATE2 к R3, после чего клиент выводит K3. Каждое рукопожатие создаёт отдельные ключи между клиентом и целевым узлом; уже задействованные узлы пересылают сообщения для следующих переходов, обработав собственный слой. Пунктирные стрелки изображают логический путь через промежуточные узлы, а не прямое соединение. Между соседними узлами действуют отдельные TLS-каналы.}
% Телескопическое построение цепи.
\begin{tikzpicture}[font=\sffamily\small,>=Stealth,
  msg/.style={font=\ttfamily\scriptsize,fill=white,inner sep=1pt},
  indietro/.style={-{Stealth[length=2.5mm]},thick,torverde!60!black},
  circav/.style={cellaf,dashed},
  circind/.style={indietro,dashed}]
  \foreach \n/\x/\t in {c/0/Клиент,r1/3.4/$R_1$ сторожевой,r2/6.8/$R_2$ промежуточный,r3/10.2/$R_3$ выходной}
    \node[font=\sffamily\small] (\n) at (\x,0) {\t};
  \foreach \x in {0,3.4,6.8,10.2} \draw[torgrigio] (\x,-0.3) -- (\x,-5.5);
  % первый переход
  \draw[cellaf] (0,-0.9) -- node[msg,above]{CREATE2} (3.4,-0.9);
  \draw[indietro] (3.4,-1.5) -- node[msg,above]{CREATED2} (0,-1.5);
  \node[anchor=east,font=\scriptsize,text=torgrigio] at (-0.2,-1.2) {ключи $K_1$};
  % второй переход
  \draw[circav] (0,-2.5) -- node[msg,above]{EXTEND2} (3.4,-2.5);
  \draw[cellaf] (3.4,-2.5) -- node[msg,above]{CREATE2} (6.8,-2.5);
  \draw[indietro] (6.8,-3.1) -- node[msg,above]{CREATED2} (3.4,-3.1);
  \draw[circind] (3.4,-3.1) -- node[msg,above]{EXTENDED2} (0,-3.1);
  \node[anchor=east,font=\scriptsize,text=torgrigio] at (-0.2,-2.8) {ключи $K_2$};
  % третий переход
  \draw[circav] (0,-4.2) -- node[msg,above]{EXTEND2} (6.8,-4.2);
  \draw[cellaf] (6.8,-4.2) -- node[msg,above]{CREATE2} (10.2,-4.2);
  \draw[indietro] (10.2,-4.8) -- node[msg,above]{CREATED2} (6.8,-4.8);
  \draw[circind] (6.8,-4.8) -- node[msg,above]{EXTENDED2} (0,-4.8);
  \node[anchor=east,font=\scriptsize,text=torgrigio] at (-0.2,-4.5) {ключи $K_3$};
  % каналы
  \foreach \a/\b/\y in {0/3.4/-5.8,3.4/6.8/-5.8,6.8/10.2/-5.8}
    \draw[thick,torgrigio!60,{Bar}-{Bar}] (\a,\y) -- node[below,font=\scriptsize]{TLS-канал} (\b,\y);
  \node[font=\scriptsize,text=torgrigio,align=center] at (5.1,-6.6) {каждое рукопожатие — между клиентом и одним узлом};
\end{tikzpicture}
